\section{Uniform Riccati and boundary estimates}
\label{sec:uniform}

The finite premises now have to control arbitrarily long interiors.
Two estimates are needed: contraction of the Riccati map keeps every
pivot positive, and decay of the endpoint responses makes the accumulated
boundary corrections summable. The column-weight estimate gives the
first, and the row-weight estimate gives the second. The same analytic
argument applies separately to the two parameter columns above.
For a real symmetric matrix $U$, use the order-unit norm
\[
 |U|_P=\inf\{d\ge0:-dP\preceq U\preceq dP\}
      =\norm{P^{-1/2}UP^{-1/2}}_2.
\]
On the closed ball $\mathcal K=\{X=X^T:|X-Z|_P\le\delta\}$,
$P\prec2I$ gives $\norm{X-Z}_2\le2\delta=:e$.
The center bounds imply $X\succ I/3$, $\norm{X^{-1}}_2\le3$, and
\[
 \norm{X^{-1}-Z^{-1}}_2\le6e.
\]
Since $\norm{E_\pm}_F^2=14$, putting $Y_X=F_+(X)$ yields
\[
 \norm{Y_X-Y}_2\le84e<1/6,\qquad
 Y_X\succ I/3,\qquad
 \norm{Y_X^{-1}-Y^{-1}}_2\le504e.
\]
For $T(X)=X^{-1}E_+Y_X^{-1}E_-$, a two-term product difference gives
\begin{equation}
 \norm{T(X)-T_Z}_2\le14(6\cdot3+2\cdot504)e=14364e.
 \label{eq:Tperturb}
\end{equation}
The norm conversion factor from the Euclidean norm to either
\[
 \norm T_{P,\mathrm{col}}=\norm{P^{1/2}TP^{-1/2}}_2,
 \qquad
 \norm T_{Q,\mathrm{row}}=\norm{Q^{-1/2}TQ^{1/2}}_2
\]
is at most $\sqrt{20/9}<3/2$. Thus the perturbation in either norm is
less than $43092\delta<10^{-4}$. Both parameter columns satisfy
$\sqrt\kappa+10^{-4}<q$, and hence throughout the ball
\begin{equation}
 T(X)^TPT(X)\prec q^2P,\qquad T(X)QT(X)^T\prec q^2Q.
 \label{eq:dual-contractions}
\end{equation}

The derivative is $D\Phi_X[U]=T(X)^TUT(X)$. The first inequality in
\eqref{eq:dual-contractions} therefore bounds its order-unit norm by
$\theta|U|_P$, where $\theta=q^2$. Integration along a line segment in
$\mathcal K$ proves the same Lipschitz bound for $\Phi$. The center
residual has order-unit norm below $\delta/36$. Both columns satisfy
$1/36+\theta<1$, so
\[
 |\Phi(X)-Z|_P<\delta/36+\theta\delta<\delta.
\]
Banach's theorem gives a fixed point $X_*\in\mathcal K$. The actual
even orbit starts at $Z$, whence
\begin{equation}
 \norm{X_{J+2s}-X_*}_2\le4\delta\theta^s\qquad(s\ge0).
 \label{eq:pivot-decay}
\end{equation}
All these pivots and their intermediate odd partners exceed $I/3$.
Together with Lemma~\ref{lem:finite-premises}, this proves positivity of
every pivot at either energy.

\subsection{Response decay and the limiting core}

Two steps in \eqref{eq:recurrence} send $R$ to $RT(X)$ and $W$ to
$T(X)^TW$. Their natural norms are
\[
 \norm{RQ^{1/2}}_2,\qquad \norm{Q^{1/2}W}_2.
\]
The second inequality in \eqref{eq:dual-contractions} contracts both
by $q$. The table satisfies $a^2>(10/9)\zeta$; the response premises and
$Q\succ(9/10)I$ give
\begin{equation}
 \norm{R_{J+2s}}_2,\ \norm{W_{J+2s}}_2<a q^s.
 \label{eq:even-response}
\end{equation}
A one-step transfer has norm at most $3\sqrt{14}<12$. Consequently
\begin{equation}
 \norm{R_a}_2,\ \norm{W_a}_2\le bq^s
 \quad\text{for }a\in\{J+2s,J+2s+1\}.
 \label{eq:pair-response}
\end{equation}
The orientation $TQT^T$ is essential; a column-only estimate would not
establish these two response bounds.

Define absolutely convergent series along the actual open trajectory:
\begin{align*}
 G_\infty&=G_0-\sum_{a\ge0}R_aX_a^{-1}R_a^T,\\
 H_\infty&=H_0-\sum_{a\ge0}W_a^TX_a^{-1}W_a,\\
 C_\infty&=C_0-\sum_{a\ge0}R_aX_a^{-1}W_a.
\end{align*}
Convergence follows from \eqref{eq:pair-response} and
$\norm{X_a^{-1}}_2\le3$ after the entrance. Set
\begin{equation}
 S_\infty=\begin{pmatrix}
 G_\infty&C_\infty\\C_\infty^T&H_\infty-E_+^TX_*^{-1}E_+
 \end{pmatrix}.
 \label{eq:Sinfty}
\end{equation}
There is no replacement of every pivot in the sums by $X_*$.
Common initial terms therefore cancel exactly when a finite core is
compared with this limit.

\begin{proposition}[Complete one-cell tail]\label{prop:tail}
For even $\ell\ge J+2$, write $p=\ell-2=J+2s$. Then
\begin{align}
 \norm{S_\ell(c)-S_\infty}_2
 &\le B_s,\notag\\
 B_s&=\frac{24b^2q^{2s}}{1-q^2}+48aq^s+576\delta\theta^s
      <\epsilon.\label{eq:Btail}
\end{align}
Furthermore $S_\infty\succ(\gamma-\epsilon)I_6$ and
$S_\ell(c)\succ(\gamma-2\epsilon)I_6$.
\end{proposition}
\begin{proof}
Each quadratic or mixed Schur suffix after $p$ is bounded by
\[
 T_s=\frac{6b^2q^{2s}}{1-q^2}.
\]
There are two single-block increments per pair, each at most
$3b^2q^{2v}$; including the term at $p$ only enlarges the bound.
The terminal corrections in \eqref{eq:single-core} have norms at most
\[
 \norm{R_pX_p^{-1}E_+}_2\le12a q^s,
 \qquad
 \norm{W_p^TX_p^{-1}E_++E_+^TX_p^{-1}W_p}_2\le24a q^s.
\]
The inverse identity and \eqref{eq:pivot-decay} give
\[
 \norm{X_p^{-1}-X_*^{-1}}_2\le36\delta\theta^s,
 \qquad
 \norm{E_+^T(X_p^{-1}-X_*^{-1})E_+}_2\le576\delta\theta^s.
\]
Thus the errors in the $G,C,H$ blocks are bounded respectively by
$T_s$, $T_s+12aq^s$, and $T_s+24aq^s+576\delta\theta^s$.
Using
\[
 \norm{\begin{pmatrix}U&V\\V^*&Z\end{pmatrix}}_2
 \le\norm U_2+2\norm V_2+\norm Z_2
\]
proves \eqref{eq:Btail}. Its right side decreases with $s$, and exact
rational arithmetic gives
\begin{equation}
 B_0=
 \begin{cases}
 251089/156250000<1/500,&c=\czero,\\
 583333407/109375000000000000<10^{-8},&c=\cone.
 \end{cases}
 \label{eq:B0}
\end{equation}
Compare the certified seed $S_{J+2}(c)\succ\gamma I_6$ to the limit,
and then the limit to any other finite core. These are two separate
errors, giving the two stated lower bounds.
\end{proof}

The resulting core margins are $2/125$ and $49/50000000$, respectively.
Schur congruence transfers positivity to the full matrix; it does not
transfer these numbers as Euclidean eigenvalue gaps of the full matrix.

\subsection{Every unit complex phase}
\label{sec:phase}

For $h=8j+2$ and $|z|=1$, let $A_h(z)$ be the Hermitian cell operator
induced by
\[
 (Au)_i=u_{i-1}+u_{i+1}+\tau_{i-2}u_{i-2}+\tau_i u_{i+2},
 \qquad u_{i+h}=zu_i,\qquad \tau=w_j\text{ periodically extended}.
\]
At $h\ge18$, the two wrap blocks of $cI-A_h(z)^2$ are
$C_0(z)=\bar z C_0$ and $W_0(z)=\bar z W_0$; all other initial blocks
are unchanged. This follows directly from the two-walk paths, each wrap
path crossing the cut once in the negative direction.
The Hermitian Schur recurrence gives
\[
 X_p(z)=X_p,\ R_p(z)=R_p,\ G_p(z)=G_p,\ H_p(z)=H_p,
 \quad C_p(z)=\bar z C_p,\ W_p(z)=\bar z W_p.
\]
Invariance of $H_p$ uses conjugate transpose and $|z|=1$.
The physical terminal coupling remains $E=E_+$, so it is added to
$\bar zW_p$, not multiplied by the phase. Put
\begin{align*}
 g_p&=G_p-R_pX_p^{-1}R_p^T,&
 c_p&=C_p-R_pX_p^{-1}W_p,\\
 h_p&=H_p-W_p^TX_p^{-1}W_p-E^TX_p^{-1}E,&
 u_p&=R_pX_p^{-1}E,\qquad v_p=W_p^TX_p^{-1}E.
\end{align*}
The exact terminal core and its unitary conjugate are
\begin{align}
 S_\ell(z)&=\begin{pmatrix}
 g_p&\bar z c_p-u_p\\ *&h_p-zv_p-\bar z v_p^T
 \end{pmatrix},\notag\\
 U_z^*S_\ell(z)U_z&=\begin{pmatrix}
 g_p&c_p-zu_p\\ *&h_p-zv_p-\bar z v_p^T
 \end{pmatrix},\quad U_z=\diag(I_2,zI_4).
 \label{eq:phase-core}
\end{align}
The lower-left blocks are conjugate transposes. The limit of
\eqref{eq:phase-core} is the same $S_\infty$ for every phase. Moreover,
all the finite tail estimates in Proposition~\ref{prop:tail} are unchanged
by the factors $z$, giving a uniform bound before taking any limit.

\begin{corollary}\label{cor:phase}
For every $|z|=1$,
\[
 h\ge106\Longrightarrow\rho(A_h(z))^2<\czero,
 \qquad h\ge202\Longrightarrow\rho(A_h(z))^2<\cone,
\]
where $h\equiv2\pmod8$ and the word is $w_{(h-2)/8}$.
For any number of identical cells and either holonomy, the corresponding
finite signing obeys the same applicable bound.
\end{corollary}
\begin{proof}
The uniformly positive pivots and phase core imply $cI-A_h(z)^2\succ0$.
For $r$ identical cells with total holonomy $\alpha$, the maps
\[
 (F_zv)_{ah+s}=r^{-1/2}z^a v_s\qquad(z^r=\alpha)
\]
give an orthogonal decomposition of the full graph into the $A_h(z)$
fibers, by the same finite geometric-sum argument as
Lemma~\ref{lem:bloch}. The seam identity is precisely $z^r=\alpha$.
Every fiber is bounded, so the full radius is bounded.
\end{proof}
